% Spis skrótów (opcjonalny) - włącz przełącznikiem \showabbreviationstrue w config/metadata.tex.
% Wpisy układaj alfabetycznie. Tytuł spisu pochodzi z \abbrevname (preamble/polish.tex).
\UnnumberedChapter{\abbrevname}

\begin{abbreviations}
    \item[AMSI]    Antimalware Scan Interface (interfejs skanowania antywirusowego systemu Windows)
    \item[C2]      Command and Control (infrastruktura dowodzenia i kontroli)
    \item[DFIR]    Digital Forensics and Incident Response (informatyka śledcza i reagowanie na incydenty)
    \item[DLL]     Dynamic-Link Library (biblioteka łączona dynamicznie)
    \item[DNS]     Domain Name System (system nazw domenowych)
    \item[HTTP]    Hypertext Transfer Protocol (protokół przesyłania hipertekstu)
    \item[HTTPS]   Hypertext Transfer Protocol Secure (szyfrowany protokół przesyłania hipertekstu)
    \item[IoA]     Indicator of Activity (wskaźnik aktywności)
    \item[IoC]     Indicator of Compromise (wskaźnik kompromitacji)
    \item[IP]      Internet Protocol (protokół internetowy)
    \item[LOLBins] Living off the Land Binaries (systemowe pliki wykonywalne wykorzystywane w atakach)
    \item[PE]      Portable Executable (format plików wykonywalnych systemu Windows)
    \item[RAM]     Random-Access Memory (pamięć operacyjna o dostępie swobodnym)
    \item[SOC]     Security Operations Center (centrum operacji bezpieczeństwa)
    \item[TCP]     Transmission Control Protocol (protokół kontroli transmisji)
    \item[TTP]     Tactics, Techniques, and Procedures (taktyki, techniki i procedury)
    \item[WAN]     Wide Area Network (sieć rozległa)
\end{abbreviations}
